\chapter{Long-sightedness}
\index{Long-sightedness}

\section*{Definition and Overview}
Long-sightedness, clinically known as hyperopia, is a common refractive error of the eye where distant objects can typically be seen clearly, but near objects appear blurred. This condition occurs when the eyeball is too short from front to back, or the cornea has insufficient curvature, causing incoming light rays to focus behind the retina rather than directly on its surface. While mild hyperopia is often compensated for by the eye's natural focusing mechanism (accommodation), higher degrees of hyperopia or age-related changes can lead to persistent visual discomfort, eye strain, and headaches, necessitating optical correction.

\section*{Epidemiology}
Hyperopia is highly prevalent across all age groups globally. It is particularly common in infants and young children, where it is often a normal developmental stage that resolves as the eyeball grows (a process known as emmetropisation). In many countries, refractive errors are the most common cause of visual impairment, with hyperopia affecting a significant proportion of the population, especially older adults. The prevalence of symptomatic long-sightedness increases with age as the natural accommodating power of the lens diminishes, a process that eventually merges with presbyopia in middle-aged individuals.

\section*{Aetiology and Pathophysiology}
Hyperopia arises from structural variations in the eye's optical components, preventing parallel light rays from focusing precisely on the retina.
\begin{itemize}
    \item \textbf{Axial length shortening:} The eyeball is shorter than normal from front to back, causing the focal point of the lens and cornea to lie behind the retina.
    \item \textbf{Corneal flatness:} The cornea possesses insufficient curvature, resulting in inadequate refractive power to bend incoming light rays sharply enough to focus on the retina.
    \item \textbf{Genetic inheritance:} A strong family history of hyperopia is a major risk factor, as eyeball shape and size are highly heritable traits.
    \item \textbf{Age-related accommodative decline:} The ciliary muscles and crystalline lens gradually lose their ability to change shape (accommodate) to increase refractive power, exposing latent hyperopia.
    \item \textbf{Systemic conditions:} Rarely, hyperopia can be associated with underlying developmental anomalies or systemic syndromes, such as microphthalmia.
\end{itemize}

\section*{Clinical Features}
Symptoms of long-sightedness vary depending on the severity of the refractive error and the individual's age and accommodative capacity.
\begin{itemize}
    \item \textbf{Blurred near vision:} Difficulty focusing on close tasks, such as reading, sewing, or using a computer, while distant vision remains clear.
    \item \textbf{Asthenopia (eye strain):} A dull ache, fatigue, or burning sensation in the eyes, particularly after prolonged periods of near work.
    \item \textbf{Headaches:} Frontal or temporal headaches provoked by tasks requiring sustained visual focus and accommodation.
    \item \textbf{Squinting:} Involuntary narrowing of the eyelids to reduce the size of the blur circle and improve focus on near objects.
    \item \textbf{Amblyopia and strabismus:} In children, high uncorrected hyperopia can lead to inward turning of the eyes (accommodative esotropia) or lazy eye (amblyopia) due to constant defocus.
\end{itemize}

\section*{Differential Diagnosis}
Hyperopia must be distinguished from other refractive errors and age-related ocular changes.
\begin{itemize}
    \item \textbf{Presbyopia:} Age-related loss of lens flexibility that impairs near vision starting in the 40s; unlike hyperopia, it is not caused by a short eyeball and affects previously normal-sighted people.
    \item \textbf{Astigmatism:} A refractive error caused by an irregularly curved cornea or lens, which causes blurring at all distances rather than primarily at near.
    \item \textbf{Myopia (short-sightedness):} The opposite refractive error, where near vision is clear but distant vision is blurred.
    \item \textbf{Cataracts:} Clouding of the natural lens that causes progressive blurriness, glare, and reduced visual acuity that cannot be fully corrected with glasses.
    \item \textbf{Pseudomyopia:} Accommodative spasm that mimics near-focusing difficulties but represents a transient functional spasm rather than a structural refractive error.
\end{itemize}

\section*{When to Seek Medical Care}
Patients should consult an optometrist or ophthalmologist if they experience progressive blurring of vision, recurrent eye strain, or headaches during reading or near work. Children should have their vision tested if they show signs of squinting, eye rubbing, or school performance issues.

\textbf{Contact your local emergency services for:}
\begin{itemize}
    \item Sudden, painless, or painful loss of vision in one or both eyes
    \item Sudden onset of double vision or severe eye pain
    \item Persistent flashes of light or a sudden shower of floaters in the visual field
    \item Physical trauma to the eye or chemical exposure
\end{itemize}

\section*{Diagnosis and Investigations}
Diagnosis is established through a comprehensive eye examination by an eye care professional.
\begin{itemize}
    \item \textbf{Visual acuity testing:} Assessing the patient's ability to read letters on a Snellen chart at a distance and a near reading card.
    \item \textbf{Subjective refraction:} Utilising a phoropter or trial lenses to determine the exact lens power (in dioptres) required to correct the refractive error.
    \item \textbf{Objective refraction:} Performing retinoscopy or using an autorefractor to measure the eye's refractive state without patient input.
    \item \textbf{Cycloplegic refraction:} Using eye drops to temporarily paralyse the ciliary muscle, essential in children to uncover latent hyperopia.
    \item \textbf{Slit-lamp examination:} Assessing the health of the anterior segment of the eye, including the cornea, iris, and lens.
    \item \textbf{Ophthalmoscopy:} Examining the retina and optic nerve to ensure the absence of structural ocular diseases.
\end{itemize}

\section*{Management}
Management aims to redirect light rays onto the retina, restoring clear near vision and relieving symptoms of eye strain.
\begin{itemize}
    \item \textbf{Corrective spectacles:} Prescription glasses featuring convex (plus) lenses to increase the eye's refractive power and bring the focal point forward onto the retina.
    \item \textbf{Contact lenses:} Soft or rigid gas-permeable contact lenses, which provide a wider field of view and are preferred by some patients for cosmetic or lifestyle reasons.
    \item \textbf{Refractive surgery:} Laser procedures such as LASIK (laser-assisted in situ keratomileusis) or PRK (photorefractive keratectomy) to reshape the peripheral cornea and increase its curvature.
    \item \textbf{Clear lens extraction or phakic intraocular lenses:} Surgical insertion of an artificial lens, typically reserved for severe hyperopia or older adults who are not candidates for laser surgery.
    \item \textbf{Vision therapy:} In children with associated binocular vision anomalies, specialised exercises may be used alongside optical correction.
\end{itemize}

\section*{Complications}
\begin{itemize}
    \item Accommodative esotropia (inward eye turn) in children due to excessive accommodative effort.
    \item Refractive amblyopia (lazy eye) if high hyperopia is asymmetric and left uncorrected during childhood development.
    \item Increased risk of angle-closure glaucoma due to the physically smaller anterior chamber characteristic of hyperopic eyes.
    \item Chronic eye strain and tension-type headaches, impairing work or academic performance.
    \item Progressive visual impairment and difficulty performing daily activities without corrective lenses.
\end{itemize}

\section*{Prognosis and Outcomes}
The prognosis for hyperopia is excellent. Visual acuity and comfort are almost always fully restored with appropriate spectacles, contact lenses, or refractive surgery. In children, early detection and correction prevent long-term complications such as amblyopia and strabismus. Hyperopia typically remains stable during adulthood but symptoms may increase in middle age as the natural accommodating power declines, requiring adjustments to the corrective lens prescription.

\section*{Prevention and Self-Care}
\begin{itemize}
    \item \textbf{Regular eye examinations:} Schedule comprehensive eye checks every two years, or annually for children and older adults.
    \item \textbf{Proper lighting:} Ensure adequate illumination when reading or performing detailed close-up tasks to minimise eye strain.
    \item \textbf{Visual hygiene (20-20-20 rule):} Take a 20-second break every 20 minutes to look at an object 20 feet (6 metres) away to relax ciliary muscles.
    \item \textbf{Correct prescription wear:} Wear prescribed glasses or contact lenses as directed by your eye care professional.
    \item \textbf{Ergonomic work setup:} Position computer screens at an appropriate distance and angle to reduce visual fatigue.
\end{itemize}

\section*{Sources and Further Reading}
The following reputable organisations provide further information on hyperopia and eye health:
\begin{itemize}
    \item Optometry Australia. \textit{Hyperopia (long-sightedness) guidelines.} https://www.optometry.org.au/
    \item Macular Disease Foundation Australia. \textit{Resources on general eye health and refractive errors.} https://www.mdfoundation.com.au/
    \item Mayo Clinic. \textit{Farsightedness (hyperopia) overview.} https://www.mayoclinic.org/
    \item NHS. \textit{Long-sightedness causes, symptoms, and correction.} https://www.nhs.uk/
    \item American Academy of Ophthalmology. \textit{Hyperopia clinical guidelines and patient resources.} https://www.aao.org/
\end{itemize}
